\section{积分算子的谱展开}
\label{sec:ospec-integral-expansion}

第~\ref{chap:linear-operators}章证明了：$\sigma$有限测度空间上的平方可积核给出有界的紧积分算子，伴随的核是原核交换变量后取共轭，核的$L^2$误差控制算子范数误差。本节把前几节的展开落实到这些核上。算子的展开是关于输入函数$f$的级数；核是乘积空间上的一个函数，它也可以展开成乘积函数的级数。两种展开所用的收敛意义不同，这是本节要分清的主要问题。

下面设$(X,\mu)$、$(Y,\nu)$是$\sigma$有限测度空间，$k\in L^2(X\times Y,\mu\otimes\nu)$，$T:L^2(Y,\nu)\to L^2(X,\mu)$是以$k$为核的积分算子。对$a\in L^2(\mu)$、$b\in L^2(\nu)$，记$(a\otimes\bar b)(x,y)=a(x)\overline{b(y)}$。先建立两个把“核”与“算子”联系起来的工具。

\begin{lemma}[乘积函数与核系数]
\label{lem:ospec-product-coefficients}
设$(a_j)\subseteq L^2(\mu)$、$(b_l)\subseteq L^2(\nu)$都是标准正交序列。则$\{a_j\otimes\bar b_l\}$是$L^2(\mu\otimes\nu)$中的标准正交族，并且
\[
\langle a_j\otimes\bar b_l,k\rangle_{L^2(\mu\otimes\nu)}=\langle a_j,Tb_l\rangle_{L^2(\mu)} .
\]
\end{lemma}

\begin{proof}
函数$\lvert a_j(x)\rvert^2\lvert b_l(y)\rvert^2$的积分由Tonelli定理等于$\lVert a_j\rVert^2\lVert b_l\rVert^2=1$，所以$a_j\otimes\bar b_l\in L^2(\mu\otimes\nu)$。两个乘积函数之积可积（Cauchy--Schwarz），由Fubini定理分离变量，
\begin{equation*}
\begin{aligned}
\langle a_j\otimes\bar b_l,a_{j^\prime}\otimes\bar b_{l^\prime}\rangle&=\int\overline{a_j}a_{j^\prime}\,\mathrm d\mu\cdot\int b_l\overline{b_{l^\prime}}\,\mathrm d\nu\\
&=\langle a_j,a_{j^\prime}\rangle\,\langle b_{l^\prime},b_l\rangle=\delta_{j,j^\prime}\delta_{l,l^\prime}.
\end{aligned}
\end{equation*}
同样，$\overline{a_j(x)}b_l(y)k(x,y)$在乘积空间上可积，先对$y$积分得
\begin{equation*}
\int_X\overline{a_j(x)}(Tb_l)(x)\,\mathrm d\mu(x)=\langle a_j,Tb_l\rangle.
\end{equation*}
\end{proof}

\begin{lemma}[核由算子唯一确定]
\label{lem:ospec-kernel-unique}
若$r\in L^2(\mu\otimes\nu)$给出的积分算子为零，则$r=0$在$\mu\otimes\nu$几乎处处成立。
\end{lemma}

\begin{proof}
对$\mu(A)<\infty$、$\nu(B)<\infty$，有$\langle\mathbf 1_A\otimes\mathbf 1_B,r\rangle=\langle\mathbf 1_A,T_r\mathbf 1_B\rangle=0$，计算方式同上一引理。第~\ref{chap:linear-operators}章证明了这类矩形示性函数的线性组合在$L^2(\mu\otimes\nu)$中稠密，所以$r$与一个稠密集正交，只能为零。
\end{proof}

\subsection{自伴积分算子的特征展开}
\label{sec:ospec-self-adjoint-kernel}

取$X=Y$、$\mu=\nu$，并设$k(x,y)=\overline{k(y,x)}$对$\mu\otimes\mu$几乎处处成立。由第~\ref{chap:linear-operators}章，$T$是$L^2(\mu)$上的紧自伴算子，定理~\ref{thm:ospec-compact-self-adjoint}给出标准正交特征函数$(e_j)$与非零实特征值$(\lambda_j)$。

\begin{theorem}[自伴积分算子的特征展开]
\label{thm:ospec-self-adjoint-kernel}
在上述条件下，
\begin{enumerate}
\item 对每个$f\in L^2(\mu)$，$Tf=\sum_j\lambda_j\langle e_j,f\rangle e_j$，级数在$L^2(\mu)$中收敛；
\item 核满足$k=\sum_j\lambda_j\,e_j\otimes\bar e_j$，即$k(x,y)=\sum_j\lambda_je_j(x)\overline{e_j(y)}$，级数在$L^2(\mu\otimes\mu)$中收敛；
\item $\lVert k\rVert_{L^2(\mu\otimes\mu)}^2=\sum_j\lambda_j^2$，截去$j>m$各项后核的$L^2$误差为$\bigl(\sum_{j>m}\lambda_j^2\bigr)^{1/2}$。
\end{enumerate}
\end{theorem}

\begin{proof}
（1）就是式\eqref{eq:ospec-spectral-expansion}。（2）由引理~\ref{lem:ospec-product-coefficients}，$k$关于标准正交族$\{e_j\otimes\bar e_l\}$的系数为$\langle e_j,Te_l\rangle=\lambda_l\delta_{j,l}$。由Bessel不等式$\sum_j\lambda_j^2\leq\lVert k\rVert^2<\infty$，所以$\tilde k=\sum_j\lambda_je_j\otimes\bar e_j$在$L^2(\mu\otimes\mu)$中收敛。

记$k_m$为前$m$项的部分和，它给出算子$T_{k_m}f=\sum_{j\leq m}\lambda_j\langle e_j,f\rangle e_j$。由核误差控制算子误差，$T_{k_m}\to T_{\tilde k}$按算子范数收敛；由（1），$T_{k_m}f\to Tf$对每个$f$成立。所以$T_{\tilde k}=T$，核$k-\tilde k$给出零算子，由引理~\ref{lem:ospec-kernel-unique}，$k=\tilde k$几乎处处。

（3）由正交展开的勾股定理，$\lVert k\rVert^2=\lVert\tilde k\rVert^2=\sum_j\lambda_j^2$，尾部的范数平方为$\sum_{j>m}\lambda_j^2$。
\end{proof}

定理的第（3）条与上一节的截断误差合在一起，给出同一个截断的两种误差：算子范数误差$\lVert T-T_m\rVert=\lvert\lambda_{m+1}\rvert$只看第一个被舍弃的特征值，核的$L^2$误差$(\sum_{j>m}\lambda_j^2)^{1/2}$要把全部舍弃项的平方加起来，后者总不小于前者。

贯穿例子中，$1+xz=1\cdot1+\tfrac13(\sqrt3x)(\sqrt3z)$正是两项的特征展开，$\lVert k\rVert^2=10/9=1^2+(1/3)^2$。第~\ref{chap:linear-operators}章用乘积积分直接算出的$10/9$在这里得到了解释：它是特征值的平方和。截去第二项，核误差与算子误差都等于$1/3$，因为只舍弃了一项。

\paragraph{收敛意义的边界}
核的展开在$L^2(\mu\otimes\mu)$中收敛，意思是$\iint\lvert k-k_m\rvert^2\to0$。它不推出$k_m(x,y)\to k(x,y)$对每一对$(x,y)$成立，更不推出一致收敛。这不只是技术细节：平方可积核本身只是几乎处处的等价类，在零测集上改动它不改变积分算子。例如在$[-1,1]^2$上令$k^\prime(x,z)=1+xz-10\cdot\mathbf 1_{\{x=z\}}$，对角线在乘积测度下为零测集，$k^\prime$与$k$给出同一个正算子，但$k^\prime(x,x)=x^2-9<0$。所以从本节的结论出发，不能谈论核在对角线上的取值，也不能断言展开在每一点成立。

逐点或一致收敛需要附加条件，例如核连续、算子为正，并在紧区域上考虑连续的特征函数代表；这类结论属于第~\ref{chap:rkhs}章的Mercer定理，不由这里的谱分解自动得到。

\subsection{非自伴积分算子的特征展开}
\label{sec:ospec-non-self-adjoint-kernel}

自伴条件要求核关于对角线共轭对称，很多核不满足它。上一节的累积积分说明，一般非自伴紧算子可能连一个特征向量都没有，更没有完备的特征系统。本小节先给出一类仍有正交特征展开的非自伴算子，再说明其余情形改用奇异分解。

\paragraph{正规算子}
设$H$是复Hilbert空间。若$T\in\mathcal B(H)$满足$T^*T=TT^*$，则称$T$为\term{正规算子}（normal operator）。自伴算子显然正规；反自伴算子$T^*=-T$也正规，却不自伴（$T\neq0$时）。对积分算子，由伴随核公式，$T^*T$与$TT^*$的核分别是$\int\overline{k(z,x)}k(z,y)\,\mathrm d\mu(z)$与$\int k(x,z)\overline{k(y,z)}\,\mathrm d\mu(z)$，正规条件就是这两个核几乎处处相等。

\begin{theorem}[紧正规算子的谱分解]
\label{thm:ospec-compact-normal}
设$T$是复Hilbert空间$H$上的紧正规算子。则存在有限或可数的标准正交序列$(e_j)$及非零复数$\mu_j$（序列无限时$\mu_j\to0$），使$Te_j=\mu_je_j$，$(e_j)$的闭张成等于$(\ker T)^\perp$，并且$Tu=\sum_j\mu_j\langle e_j,u\rangle e_j$对每个$u\in H$成立。
\end{theorem}

\begin{proof}
令$A=(T+T^*)/2$、$B=(T-T^*)/(2\mathrm i)$。二者都是紧自伴算子，$T=A+\mathrm iB$，并且正规条件推出$AB=BA$。由$T^*T=(A-\mathrm iB)(A+\mathrm iB)=A^2+B^2$，得$\lVert Tu\rVert^2=\lVert Au\rVert^2+\lVert Bu\rVert^2$，所以$\ker T=\ker A\cap\ker B$。

对$A$的每个非零特征值$\alpha$，特征空间$E_\alpha$有限维，并且在$B$下不变：若$Au=\alpha u$，则$A(Bu)=BAu=\alpha Bu$。$B$在$E_\alpha$上的限制是有限维自伴变换，可以取$E_\alpha$的标准正交基使它对角化；这些基向量同时是$A$与$B$的特征向量，因而是$T$的特征向量。同理，$\ker A$在$B$下不变，$B$在$\ker A$上的限制是紧自伴算子，由定理~\ref{thm:ospec-compact-self-adjoint}，它的非零特征方向与零空间$\ker A\cap\ker B$一起张满$\ker A$，这些特征方向上$T=\mathrm iB$。

把两部分的特征向量合在一起，得到标准正交序列$(e_j)$，其闭张成是$\ker T=\ker A\cap\ker B$的正交补，相应特征值$\mu_j=\alpha_j+\mathrm i\beta_j\neq0$。展开式与$\mu_j\to0$的论证与自伴情形相同。
\end{proof}

对平方可积核$k$，若$T$正规，按定理~\ref{thm:ospec-self-adjoint-kernel}同样的证明（引理~\ref{lem:ospec-product-coefficients}给出系数$\langle e_j,Te_l\rangle=\mu_l\delta_{j,l}$），得到核在乘积$L^2$中的正交特征函数展开
\[
k=\sum_j\mu_j\,e_j\otimes\bar e_j,\qquad\lVert k\rVert_{L^2(\mu\otimes\mu)}^2=\sum_j\lvert\mu_j\rvert^2 .
\]

\paragraph{周期卷积核}
在$[0,1)$上取Lebesgue测度，令$h$为$t-\tfrac12$（$t\in[0,1)$）的$1$周期延拓，核$k(x,y)=h(x-y)$。对非整数$t$有$h(-t)=-h(t)$，所以$k(y,x)=-k(x,y)$几乎处处，伴随核为$-k$，$T^*=-T$：$T$正规而不自伴。

特征函数可以直接验证。令$e_n(x)=\mathrm e^{2\pi\mathrm inx}$，$n\in\mathbb Z$。由周期性作变量替换$t=x-y$，
\[
(Te_n)(x)=\int_0^1h(x-y)\mathrm e^{2\pi\mathrm iny}\,\mathrm dy=\mathrm e^{2\pi\mathrm inx}\int_0^1h(t)\mathrm e^{-2\pi\mathrm int}\,\mathrm dt=\mu_ne_n(x).
\]
直接积分得$\mu_0=0$，而$n\neq0$时$\mu_n=\int_0^1t\,\mathrm e^{-2\pi\mathrm int}\,\mathrm dt=\mathrm i/(2\pi n)$。特征值都是纯虚数，这与反自伴相符。

$(e_n)$是$L^2(0,1)$的标准正交基（第~\ref{chap:fourier-transform}章），所以$Tf=\sum_{n\neq0}\mu_n\langle e_n,f\rangle e_n$，$\ker T$由常函数组成。数值核对：$\lVert k\rVert^2=\int_0^1(t-\tfrac12)^2\,\mathrm dt=1/12$，而$\sum_{n\neq0}\lvert\mu_n\rvert^2=2\sum_{n\geq1}\frac{1}{4\pi^2n^2}=\frac{1}{2\pi^2}\cdot\frac{\pi^2}{6}=\frac{1}{12}$，二者相等。

\paragraph{一般核的奇异函数展开}
正规条件不成立时，例如累积积分，或者输入、输出空间根本不同时，回到第~\ref{sec:ospec-singular}节的奇异系统。设$T:L^2(\nu)\to L^2(\mu)$的奇异系统为$(\sigma_j,v_j,u_j)$。

\begin{theorem}[平方可积核的奇异函数展开]
\label{thm:ospec-kernel-svd}
在本节开头的设定下，
\[
k=\sum_j\sigma_j\,u_j\otimes\bar v_j,\qquad\lVert k\rVert_{L^2(\mu\otimes\nu)}^2=\sum_j\sigma_j^2 ,
\]
级数在$L^2(\mu\otimes\nu)$中收敛。记$k_m$为前$m$项之和、$T_m$为相应的截断算子，则
\[
\lVert k-k_m\rVert_{L^2(\mu\otimes\nu)}=\Bigl(\sum_{j>m}\sigma_j^2\Bigr)^{1/2},\qquad\lVert T-T_m\rVert=\sigma_{m+1} .
\]
\end{theorem}

\begin{proof}
由引理~\ref{lem:ospec-product-coefficients}，$k$关于标准正交族$\{u_j\otimes\bar v_l\}$的系数为$\langle u_j,Tv_l\rangle=\sigma_l\langle u_j,u_l\rangle=\sigma_l\delta_{j,l}$。与定理~\ref{thm:ospec-self-adjoint-kernel}的证明相同，Bessel不等式保证级数收敛，部分和给出的算子按算子范数收敛到由式\eqref{eq:ospec-svd}给出的$T$，引理~\ref{lem:ospec-kernel-unique}说明两个核相等。核误差由勾股定理得到，算子误差就是定理~\ref{thm:ospec-svd-truncation}。
\end{proof}

累积积分的核$\mathbf 1_{\{t<x\}}$满足$\lVert k\rVert^2=1/2$，与$\sum_k\sigma_k^2=\frac{4}{\pi^2}\sum_k\frac{1}{(2k-1)^2}=\frac{4}{\pi^2}\cdot\frac{\pi^2}{8}=\frac12$一致。表~\ref{tab:ospec-volterra-errors}列出截断$m$项后的两种误差。

\begin{table}[!htb]
  \centering
  \caption{累积积分截断奇异展开的两种误差。算子范数误差是第一个被舍弃的奇异值$\sigma_{m+1}$，核的$L^2$误差是全部舍弃奇异值平方和的平方根，后者总不小于前者，并且下降得更慢。}
  \label{tab:ospec-volterra-errors}
  \begin{tabular}{@{}ccc@{}}
    \toprule
    保留项数$m$ & 算子误差$\sigma_{m+1}$ & 核误差$(\sum_{j>m}\sigma_j^2)^{1/2}$\\
    \midrule
    $1$ & $0.212$ & $0.308$\\
    $2$ & $0.127$ & $0.223$\\
    $3$ & $0.091$ & $0.183$\\
    $4$ & $0.071$ & $0.159$\\
    \bottomrule
  \end{tabular}
\end{table}

两种误差度量不同的对象。算子误差回答“对最坏的单位输入，近似输出差多少”，它由下一个奇异值决定，按$1/m$下降；核误差回答“作为乘积空间上的函数，近似核差多少”，它累积所有舍去的方向，在这个例子中按$1/\sqrt m$下降。用核误差估计算子误差是安全的上界，却可能过于保守。
